\scanpage{177}
\mspage{67}
\margin{176}

Note that $C_n^{\unsure{v|\nu}}$ is an even function for $n$ even and an odd function for $n$ odd, hence we have
\[
\gamma_n=\begin{cases}
\displaystyle\frac{w_p}{C_n^{1/2\,p}(1)}\,i\pi\int_0^1 C^{1/2\,p}(x)(1-x^2)^{1/2\,p-1/2}\,dx&n\text{ odd}\\[1.5ex]
\displaystyle-\frac{2w_p}{C_n^{1/2\,p}(1)}\int_0^1\log x\ C_n^{1/2\,p}(x)(1-x^2)^{1/2\,p-1/2}\,dx&n\text{ even}
\end{cases}
\]
the result is then proved by evaluating these integrals (see [18a]).

\uline{4.5 Reduction of a singular equation to the Fredholm case}.

We now return to the algebra $\mathcal{G}$ of singular operators on $M$. For each $J\in\mathcal{G}$ and each $x\in M$ we shall define the symbol of $J$ at $x$; this will be a function on $S_x$. Let $\mathbb{H}_xJ$ be the homogeneous part of $J$ at $x$; we define $\mathbb{S}_xJ=\unsure{\mathbf{S}|\mathbb{S}}\unsure{\mathbf{H}|\mathbb{H}}_xJ$, thus $\unsure{\mathbf{S}|\mathbb{S}}_x$ is the composition of the homomorphism $\mathbb{H}_x$ and the isomorphism $\mathbb{S}$.

\uline{Definition 4.5.1.} If $J\in\mathcal{G}$ then the \uline{global symbol} of $J$ is a function on the tangent sphere bundle $S(M)$ defined by
\[
\mathbb{S}_MJ(b)=\unsure{\mathbf{S}|\mathbb{S}}_{\pi(b)}J(b).
\]
Note that $\unsure{\mathbf{S}|\mathbb{S}}_MJ=0$ implies that $\mathbb{H}_xJ=0$ for all $x\in M$ and $J\in\bigcap_{x\in M}\mathcal{F}_x=\mathcal{F}$, so $J$ is completely continuous.
